\section{The ratio objective}\label{sec:ratio}

Two classical devices remove a ratio from an objective. \citet{CharnesCooper1962}
rescale a linear fractional program into a linear one by normalising the
denominator, which is exactly the transformation that turns
\(\max p(S)/w(S)\) into the program with \(\sum_i w_i y_i = 1\) that \FS{}
solves. \citet{Dinkelbach1967} instead treats the ratio parametrically, solving
\(\max\{p(S) - \lambda w(S)\}\) and updating \(\lambda\) by Newton's method;
applied to closures, each iterate is a maximum-closure problem and therefore a
minimum cut. This project implements the second as one of its reference
backends.

Maximum-ratio closure --- maximum-density or maximum-mean closed set, depending
on the community --- is therefore classical in both readings, and the
equivalence between the fractional point of view and the level-set argument
(\(y = \int \chi^{C_t}\,dt\), each \(C_t\) a closure) is the standard device
relating fractional closure points to closures. The tutorial of this project
proves it in that form and attributes it there.

\subsection{The LP class, stated precisely}\label{sec:class}

Since the positioning turns on which linear program is being specialised to, it
is worth naming the class rather than the application. The program is
\begin{equation}\label{eq:class}
  \max\ \sum_{i \in R} p_i y_i
  \quad\text{s.t.}\quad
  \sum_{i \in R} w_i y_i = 1,
  \qquad y_i - y_j \le 0 \ \ \forall (i,j) \in A[R],
  \qquad \vect{y} \ge 0,
\end{equation}
with $\vect{w} > 0$, $\vect{p}$ of arbitrary sign and no upper bounds. Three equivalent
readings of the same class:

\begin{enumerate}[label=(\alph*),leftmargin=*]
\item \textbf{One budget row over a homogeneous difference system.} The
      constraints $y_i - y_j \le 0$ form a system of homogeneous difference
      inequalities on a DAG; a single dense equality row with positive
      coefficients is added.
\item \textbf{The base of an order cone.} The set $\{\vect{y} \ge 0 : y_i \le y_j\}$ is
      the cone of nonnegative order-preserving vectors on the poset, whose
      extreme rays are the characteristic vectors of the nonempty closed sets;
      the equality slices it into the polytope whose vertices are
      $\chi^C / w(C)$. Optimising over it therefore returns
      $\max_C p(C)/w(C)$, which is why the transformation of
      \citet{CharnesCooper1962} is exact here. The polytope is the cone analogue
      of Stanley's order polytope, whose vertices are the characteristic vectors
      of the filters of the poset \citep{Stanley1986}.
\item \textbf{Transposed incidence plus one dense row.} Each difference
      constraint has a $+1$ and a $-1$ in the columns of two vertex variables, so
      the matrix is the \emph{transpose} of a digraph incidence matrix, with one
      dense row appended and a right-hand side that is zero on the network part.
\end{enumerate}

Reading (c) matters for \cref{sec:forests} and was stated loosely in an earlier
draft. In the side-constrained network flow literature the variables are arcs
and the rows are vertices; here the variables are vertices and the rows are arcs. The
two are linear programming duals of one another --- the dual of
\eqref{eq:class} is $\min \beta$ subject to $w_i\beta + \divg_i(\vecg{\alpha}) \ge p_i$,
$\vecg{\alpha} \ge 0$, a flow-like system with one scalar --- so the algebra transports
between them, but pricing on one side corresponds to the ratio test on the other
and the roles are not interchangeable without care. This is the same duality
\citet{UnderwoodTolwinski1998} exploit when they attack the ultimate pit problem
through the dual. It also places \eqref{eq:class} in the same structural family
as isotonic regression and total variation on a graph, which are difference
systems on vertex variables too --- and that is why \citet{Yang2026} is closer
prior art than the network-simplex-with-a-side-row line, not further.

\paragraph{The one non-generic modelling choice.} Replacing the equality
\(\sum_i w_i y_i = 1\) by an inequality adds the origin to the feasible set and
returns value zero whenever every nonempty closure has negative profit. With
nonnegative profits this is immaterial. Here it is not: profits carry either
sign, a profitable vertex may require loss-making prerequisites, and the outer
decomposition calls the oracle on residual graphs whose profits deteriorate, so
the last calls typically have a negative optimal ratio and the decomposition
still needs the \emph{best nonempty} closure at each of them. The equality is
deliberate and is argued as such --- but it is a modelling choice inside a
classical transformation, not a contribution.

\paragraph{Recent work on the ratio side.} \citet{Hochbaum2024} gives an
incremental parametric cut method for ratio problems of this family; it is the
most recent point of comparison found for the oracle itself and is in
\file{literature/papers/}. It is a competitor to the Dinkelbach backend, not to
\FSC{}'s basis machinery, and it should be cited where the oracle is discussed.
